% eksperymentalne tłumaczenie na włoski

%

\section{Rozwiązywanie trójkątów} % lang-pl
\section{Risoluzione dei triangoli} % lang-it
% https://en.wikipedia.org/wiki/Solution_of_triangles

% TODO: https://en.wikipedia.org/wiki/Skinny_triangle

Rozwiązywanie trójkątów (łacińskie \emph{solutio triangulorum}) to najważniejsze rachunkowe zadanie w~trygonometrii. % lang-pl
Dane są pewne informacje o trójkącie na płaszczyźnie lub sferze, takie jak: długości boków albo kąty wewnętrzne; należy ustalić miary pozostałych. % lang-pl
Zastosowania obejmują geodezję, astronomię, nawigację i nie tylko. % lang-pl
La risoluzione dei triangoli (dal latino \emph{solutio triangulorum}) è il compito computazionale più importante in trigonometria. % lang-it
Sono dati alcuni dati su un triangolo nel piano o sulla sfera, come: lunghezze dei lati o angoli interni; occorre determinare le misure degli altri. % lang-it
Le applicazioni includono geodesia, astronomia, navigazione e non solo. % lang-it

Ograniczymy się do płaszczyzny i podamy kilka nowych zależności trygonometrycznych, które pomagają w rozwiązywaniu trójkątów. % lang-pl
Znamy już twierdzenia sinusów i kosinusów; w ogólności to drugie jest bezpieczniejsze od pierwszego (ponieważ z faktu, że $\sin \alpha = \frac 1 2$ nie wynika, czy kąt $\alpha$ jest ostry, czy rozwarty, kosinus zaś jest jednoznaczny). % lang-pl
Ci limiteremo al piano e presenteremo alcune nuove relazioni trigonometriche che aiutano nella risoluzione dei triangoli. % lang-it
Conosciamo già il teorema dei seni e dei coseni; in generale il secondo è più sicuro del primo (poiché dal fatto che $\sin \alpha = \frac 1 2$ non si deduce se l'angolo $\alpha$ è acuto o ottuso, mentre il coseno è univoco). % lang-it

W każdym z poniższych stwierdzeń mamy trójkąt o kątach $\alpha, \beta, \gamma$, bokach $a, b, c$, połowie obwodu $p = \frac 1 2 (a + b + c)$, promieniu okręgu opisanego $R$ i wpisanego $r$. % lang-pl
In ciascuna delle seguenti affermazioni consideriamo un triangolo con angoli $\alpha, \beta, \gamma$, lati $a, b, c$, semiperimetro $p = \frac 1 2 (a + b + c)$, raggio del circolo circoscritto $R$ e raggio del circolo inscritto $r$. % lang-it

\begin{proposition}
    Zachodzi % lang-pl
    Vale % lang-it
    \begin{equation}
    S = 2 R^2 \sin \alpha \sin \beta \sin \gamma.
    \end{equation}
\end{proposition}

\begin{proposition}
    Zachodzi % lang-pl
    Vale % lang-it
    \begin{equation}
        p = R (\sin \alpha + \sin \beta + \sin \theta).
    \end{equation}
\end{proposition}

\begin{proposition}
    Zachodzi % lang-pl
    Vale % lang-it
    \begin{equation}
        r = 4R \sin \frac \alpha 2 \sin \frac \beta 2 \sin \frac \gamma 2.
    \end{equation}
\end{proposition}

\begin{proposition}
[wzór Newtona] % lang-pl
[formula di Newton] % lang-it
\index{wzór!Newtona}% % lang-pl
    Zachodzi % lang-pl
\index{formula!di Newton}% % lang-it
    Vale % lang-it
    \begin{equation}
        \frac{a + b}{c} = \frac{\cos \frac 1 2 (\alpha - \beta)}{\cos \frac 1 2 (\alpha + \beta)}.
    \end{equation}
\end{proposition}

\begin{proposition}
[wzór Mollweidego] % lang-pl
[formula di Mollweide] % lang-it
\index{wzór!Mollweidego}% % lang-pl
\index{formula!di Mollweide}% % lang-it
\label{wzor_mollweidego}%
    Zachodzi % lang-pl
    Vale % lang-it
    \begin{equation}
        \frac{a - b}{c} = \frac{\sin \frac 1 2 (\alpha - \beta)}{\sin \frac 1 2 (\alpha + \beta)}.
    \end{equation}
\end{proposition}

(Nie ma zgodności co do powyższych nazw). % lang-pl
Nieco bardziej geometryczną wersję wzorów poda Izaak Newton w 1707 roku, potem Friedrich von Oppel w 1746. % lang-pl
(Non c'è consenso sui nomi sopra riportati). % lang-it
Una versione leggermente più geometrica delle formule è data da Isaac Newton nel 1707, poi da Friedrich von Oppel nel 1746. % lang-it
\index[persons]{Newton, Izaak}% % lang-pl
\index[persons]{Newton, Isaac}% % lang-it
\index[persons]{Oppel@von Oppel, Friedrich}%
Wyrażenia, których używać będziemy do końca świata i jeden dzień dłużej zawdzięczymy Thomasowi Simpsonowi z 1748 roku, oraz Karlowi Mollweidemu, który opublikuje to samo w 1808 roku bez cytowania poprzedników. % lang-pl
Le espressioni che useremo fino alla fine del mondo e un giorno di più si devono a Thomas Simpson del 1748 e a Karl Mollweide, che pubblicò lo stesso nel 1808 senza citare i predecessori. % lang-it
\index[persons]{Simpson, Thomas}%
\index[persons]{Mollweide, Karl}%

\begin{proposition}
[wzór Regiomontanusa] % lang-pl
[formula di Regiomontano] % lang-it
\index{wzór!Regiomontanusa}% % lang-pl
    Zachodzi % lang-pl
\index{formula!di Regiomontano}% % lang-it
    Vale % lang-it
    \begin{equation}
        \frac{a + b}{a- b} = \frac{\tan \frac 1 2 (\alpha + \beta)}{\tan \frac 1 2 (\alpha - \beta)}.
    \end{equation}
\end{proposition}
% TODO: https://en.wikipedia.org/wiki/Law_of_tangents

Regiomontanus będzie znany gorzej jako Johannes Müller von Königsberg. % lang-pl
Regiomontano è meno noto come Johannes Müller von Königsberg. % lang-it
\index[persons]{Königsberg@von Königsberg, Johannes|see{Regiomontanus}}%
\index[persons]{Regiomontanus}%

\begin{proposition}
    Niech $\alpha, \beta, \gamma$ będą miarami kątów trójkąta o bokach $a, b, c$, polu $S$ i promieniu okręgu opisanego $R$. % lang-pl
    Wtedy % lang-pl
    Siano $\alpha, \beta, \gamma$ le misure degli angoli di un triangolo con lati $a, b, c$, area $S$ e raggio del circolo circoscritto $R$. % lang-it
    Allora % lang-it
    \begin{equation}
        \tan \alpha + \tan \beta + \tan \gamma = \frac{4S}{a^2 + b^2 + c^2 - 8R^2}
    \end{equation}
\end{proposition}

%